\documentclass[11pt]{article}
\usepackage[a4paper,margin=27mm]{geometry}
\usepackage{amsmath,amssymb,mathtools}
\usepackage{fontspec,unicode-math}
\setmainfont{Latin Modern Roman}
\setmathfont{Latin Modern Math}
\usepackage{xurl}
\usepackage[hidelinks]{hyperref}
\newcommand{\sref}[2]{\hyperlink{source:#1:#2}{[S#2]}}
\newcommand{\cataloguescope}[1]{\paragraph{Mathematical question.}#1}
\hypersetup{pdftitle={Problem statement}}
\begin{document}
\section{Does every prime number appear in the Euclid–Mullin sequence}\label{statement}

\subsection{Definitions and mathematical statement}

\paragraph{Shared notation.}$\mathbb N=\{1,2,\ldots\}$, and $\mathbb Z,\mathbb Q,\mathbb R,\mathbb C$ denote the integers, rationals, reals and complex numbers. Any other conventions are specified locally.


Define the Euclid--Mullin sequence recursively by $a_1=2$ and, for $n\ge1$, let $a_{n+1}$ be the least prime divisor of $1+\prod_{j=1}^n a_j$. The product plus one exceeds one, so the next term exists. No prime already in the product can divide this number; thus its terms are pairwise distinct.
\paragraph{Statement recorded in the supplied source.}
For every prime $p$, does there exist an index $n\ge1$ with $a_n=p$? \sref{statement}{2}
\subsection{Short English statement}
Does the recursive least-prime-factor Euclid sequence eventually contain every prime?
\subsection{Sources}
\begin{description}
\item[\hypertarget{source:statement:1}{S1}] ULAM.AI and the UnsolvedMath Contributors, UnsolvedMath: A Curated Collection of Open Mathematics Problems, catalogue entry. CC BY 4.0. \url{https://huggingface.co/datasets/ulamai/UnsolvedMath}.
\item[\hypertarget{source:statement:2}{S2}] Andrew R. Booker and Carl Pomerance, \emph{Squarefree smooth numbers and Euclidean prime generators}, arXiv:1607.01557v4 (2017) Introduction, p. 1. \url{https://arxiv.org/pdf/1607.01557}.
\item[\hypertarget{source:statement:3}{S3}] Wikipedia contributors, \emph{List of unsolved problems in mathematics}, revision 1366636767 (2026) Number theory, Does every prime number appear in the Euclid–Mullin sequence. \url{https://en.wikipedia.org/w/index.php?title=List_of_unsolved_problems_in_mathematics&oldid=1366636767#Number_theory}.
\item[\hypertarget{source:statement:4}{S4}] OEIS Foundation and contributors, \emph{The On-Line Encyclopedia of Integer Sequences}, A000945: Euclid-Mullin sequence: a(1) = 2, a(n+1) is smallest prime factor of 1 + Product\_\{k=1..n\} a(k)., official data snapshot \texttt{693b5499c9d2} (2026) Name, definition and comments. \url{https://raw.githubusercontent.com/oeis/oeisdata/693b5499c9d248f570dd2d979e8bd1d7eecf689b/seq/A000/A000945.seq}.
\end{description}
\label{end:statement}
\end{document}
